\chapter{Programs — Counting Sort}
%%% generated by the blueprint pipeline from TCSlib.Programs.Util.CountingSort
\section{Overview}

Declarations of the Lean module \texttt{TCSlib.Programs.Util.CountingSort}.

\section{Declarations}

\begin{definition}[Count of matching indices]
\lean{Light.cntEq}
\label{Light.cntEq}
\leanok
For a function from natural numbers to natural numbers, a target value, and a
natural-number bound, the count is the number of indices strictly below the bound at
which the function equals the target value.
\end{definition}

\begin{definition}[Count of smaller values]
\lean{Light.cntLt}
\label{Light.cntLt}
\leanok
For a function $k:\bbn\to\bbn$ and natural numbers $t$ and $j$, let $c(t,j)$ be the number
of indices $i<j$ for which $k(i)<t$.
\end{definition}

\begin{definition}[Stable sorting position]
\lean{Light.sortPos}
\label{Light.sortPos}
\leanok
\uses{Light.cntEq, Light.cntLt}
For a function $k:\bbn\to\bbn$ and natural numbers $N$ and $j$, define the position
associated to $j$ as the number of indices $i<N$ with $k(i)<k(j)$, plus the number of
indices $i<j$ with $k(i)=k(j)$.
\end{definition}

\begin{theorem}[Equality count at zero]
\lean{Light.cntEq_zero}
\label{Light.cntEq_zero}
\leanok
\uses{Light.cntEq}
\difficulty{1}
Let $kf:\mathbb{N}\to\mathbb{N}$ and $t\in\mathbb{N}$. There are no indices below $0$ at
which $kf$ equals $t$.
\end{theorem}

\begin{theorem}[Count over an empty range]
\lean{Light.cntLt_zero_right}
\label{Light.cntLt_zero_right}
\leanok
\uses{Light.cntLt}
\difficulty{1}
Let $kf:\mathbb{N}\to\mathbb{N}$ and $t\in\mathbb{N}$. The number of indices $i<0$
satisfying $kf(i)<t$ is zero.
\end{theorem}

\begin{theorem}[Count below zero]
\lean{Light.cntLt_zero_left}
\label{Light.cntLt_zero_left}
\leanok
\uses{Light.cntLt}
\difficulty{1}
Let $kf:\mathbb{N}\to\mathbb{N}$ and $j\in\mathbb{N}$. No index $i<j$ satisfies
$kf(i)<0$, so the count is zero.
\end{theorem}

\begin{theorem}[Successor formula for an equality count]
\lean{Light.cntEq_succ}
\label{Light.cntEq_succ}
\leanok
\uses{Light.cntEq}
\difficulty{1}
Let $kf:\mathbb{N}\to\mathbb{N}$ and $t,j\in\mathbb{N}$. The number of indices below
$j+1$ at which $kf$ equals $t$ is the number below $j$, plus $1$ if $kf(j)=t$ and plus
$0$ otherwise.
\end{theorem}

\begin{theorem}[Count after adding an index]
\lean{Light.cntLt_succ_right}
\label{Light.cntLt_succ_right}
\leanok
\uses{Light.cntLt}
\difficulty{1}
Let $kf:\mathbb{N}\to\mathbb{N}$ and $t,j\in\mathbb{N}$. The number of indices $i<j+1$
with $kf(i)<t$ equals the number with $i<j$ satisfying that inequality, plus $1$ if
$kf(j)<t$ and $0$ otherwise.
\end{theorem}

\begin{theorem}[Counting below a successor threshold]
\lean{Light.cntLt_succ_left}
\label{Light.cntLt_succ_left}
\leanok
\uses{Light.cntEq, Light.cntEq_succ, Light.cntLt, Light.cntLt_succ_right}
\difficulty{4}
Let $kf:\mathbb{N}\to\mathbb{N}$ and $t,j\in\mathbb{N}$. The number of indices $i<j$
with $kf(i)<t+1$ equals the number with $kf(i)<t$ plus the number with $kf(i)=t$.
\end{theorem}

\begin{theorem}[Monotonicity of a threshold count]
\lean{Light.cntLt_mono_left}
\label{Light.cntLt_mono_left}
\leanok
\uses{Light.cntLt}
\difficulty{1}
Let $kf:\mathbb{N}\to\mathbb{N}$ and $s,t,j\in\mathbb{N}$ with $s\leq t$. The number of
indices below $j$ at which $kf$ is less than $s$ is at most the number at which $kf$ is
less than $t$.
\end{theorem}

\begin{theorem}[Bound on a threshold count]
\lean{Light.cntLt_le}
\label{Light.cntLt_le}
\leanok
\uses{Light.cntLt}
\difficulty{1}
Let $kf:\mathbb{N}\to\mathbb{N}$ and $t,j\in\mathbb{N}$. The number of indices below $j$
at which $kf$ is less than $t$ is at most $j$.
\end{theorem}

\begin{theorem}[Prefix count under pointwise agreement]
\lean{Light.cntLt_congr}
\label{Light.cntLt_congr}
\leanok
\uses{Light.cntLt}
\difficulty{3}
Let $kf,kf':\mathbb{N}\to\mathbb{N}$, and let $s,m\in\mathbb{N}$. If $kf(i)=kf'(i)$ for
every $i<m$, then the numbers of indices below $m$ at which these functions,
respectively, are less than $s$ are equal.
\end{theorem}

\begin{theorem}[Full count below a threshold]
\lean{Light.cntLt_of_forall_lt}
\label{Light.cntLt_of_forall_lt}
\leanok
\uses{Light.cntLt}
\difficulty{1}
Let $kf:\mathbb{N}\to\mathbb{N}$ and let $N,K\in\mathbb{N}$. If $kf(i)<K$ for every
$i<N$, then exactly $N$ indices among $0,\ldots,N-1$ satisfy $kf(i)<K$.
\end{theorem}

\begin{theorem}[Strict increase in a prefix count]
\lean{Light.cntEq_lt_of_lt}
\label{Light.cntEq_lt_of_lt}
\leanok
\uses{Light.cntEq}
\difficulty{1}
Let $kf:\mathbb{N}\to\mathbb{N}$ and let $j,N\in\mathbb{N}$ with $j<N$. The number of
indices below $N$ at which $kf$ equals $kf(j)$ is strictly greater than the number of
such indices below $j$.
\end{theorem}

\begin{theorem}[Sorting position bounded by a value count]
\lean{Light.sortPos_lt_cntLt_succ}
\label{Light.sortPos_lt_cntLt_succ}
\leanok
\uses{Light.cntEq_lt_of_lt, Light.cntLt, Light.cntLt_succ_left, Light.sortPos}
\difficulty{3}
Let $k:\mathbb{N}\to\mathbb{N}$ and let $N,j\in\mathbb{N}$ with $j<N$. The position
assigned to $j$ when the first $N$ indices are sorted by their $k$-values, breaking ties
by increasing index, is strictly less than the number of indices among the first $N$
whose $k$-value is less than $k(j)+1$.
\end{theorem}

\begin{theorem}[Sorting position lies within the list]
\lean{Light.sortPos_lt}
\label{Light.sortPos_lt}
\leanok
\uses{Light.cntLt, Light.cntLt_le, Light.sortPos, Light.sortPos_lt_cntLt_succ}
\difficulty{3}
Let $k:\mathbb{N}\to\mathbb{N}$, and let $N,j\in\mathbb{N}$ with $j<N$. Rank the first
$N$ indices by their values under $k$, breaking ties by increasing index; then the
position assigned to $j$ is less than $N$.
\end{theorem}

\begin{theorem}[Ordering of distinct sorting positions]
\lean{Light.sortPos_lt_sortPos}
\label{Light.sortPos_lt_sortPos}
\leanok
\uses{Light.cntEq_lt_of_lt, Light.cntLt, Light.cntLt_mono_left, Light.sortPos, Light.sortPos_lt_cntLt_succ}
\difficulty{4}
Let $k:\mathbb{N}\to\mathbb{N}$ and $N,i,j\in\mathbb{N}$, with $i<N$. If either
$k(i)<k(j)$, or $k(i)=k(j)$ and $i<j$, then the position of $i$ in the ordering of the
first $N$ indices by increasing $k$-value, with ties broken by increasing index, is
strictly less than the position of $j$.
\end{theorem}

\begin{theorem}[Characterization of sorting order]
\lean{Light.sortPos_lt_sortPos_iff}
\label{Light.sortPos_lt_sortPos_iff}
\leanok
\uses{Light.sortPos, Light.sortPos_lt_sortPos}
\difficulty{5}
Let $k:\mathbb{N}\to\mathbb{N}$, and let $N,i,j\in\mathbb{N}$ with $i<N$ and $j<N$. The
position of $i$ in the ordering of the first $N$ indices by increasing $k$-value, with
ties broken by increasing index, is strictly less than the position of $j$ if and only
if either $k(i)<k(j)$, or $k(i)=k(j)$ and $i<j$.
\end{theorem}

\begin{theorem}[Injectivity of the sorting position]
\lean{Light.eq_of_sortPos_eq}
\label{Light.eq_of_sortPos_eq}
\leanok
\uses{Light.sortPos, Light.sortPos_lt_sortPos_iff}
\difficulty{3}
Let $kf:\mathbb N\to\mathbb N$ and $N\in\mathbb N$. For $j<N$, define its sorting
position to be the number of indices $m<N$ with $kf(m)<kf(j)$, plus the number of
indices $m<j$ with $kf(m)=kf(j)$. If $i,j<N$ have the same sorting position, then $i=j$.
\end{theorem}

\begin{theorem}[Every rank occurs as a sorting position]
\lean{Light.exists_sortPos_eq}
\label{Light.exists_sortPos_eq}
\leanok
\uses{Light.eq_of_sortPos_eq, Light.sortPos, Light.sortPos_lt}
\difficulty{4}
Let $kf:\mathbb N\to\mathbb N$ and let $N,q\in\mathbb N$ with $q<N$. Define the sorting
position of $j$ as the number of indices $m<N$ with $kf(m)<kf(j)$, plus the number of
indices $m<j$ with $kf(m)=kf(j)$. Then some $j<N$ has sorting position $q$.
\end{theorem}
